\scanpage{68}
\mspage{II.16}
% archive sheet number: 67
% lines 1–9
\deleted{groupe \unsure{à}} \inserted{$\unsure{\mathbb Z_\ell|Z_\ell}$ module $M\simeq$} $(\deleted{\unsure{\mathbb Q_\ell}}\mathbb Z_\ell)^r$ \uline{sur lequel} $I$ \uline{opère trivialement}, par un \deleted{groupe} $N$ abélien de rang $2g-r$, \inserted{can. isom. à $T_\ell(G(\unsure{S|s}))$} Comme l'action de $I$ est semi-simple, il s'ensuit que l'on a un ss-module $M\simeq\mathbb Z_\ell^r\subset T_\ell(A(\bar K))$ \uline{sur lequel} $I$ \uline{opère trivialement}, \deleted{\unsure{donc}} \inserted{avec $M\cap N=0$. On en déduit} un sous-groupe de la forme $(\unsure{\mathbb Q_\ell|Q_\ell}/\mathbb Z_\ell)^r$ du grp $A(K)$ : Par définition du modèle de Néron, il sera contenu dans $G(S)$, i.e. $M\subset N$, donc $M=0$, donc $r=0$, ce qui signifie que la réduction est \unsure{st.R|b.R}, cqfd.

\par\noindent\rule{.45\linewidth}{.4pt}\par\noindent\rule{.42\linewidth}{.4pt}\par

% lines 10–14
\uline{Remarque de Serre} Si on savait montrer que (dans le cas de la car. nulle) action résoluble implique action ess. abélienne, \inserted{(VA sur corps de type fini)} on pourrait en conclure qu'une \inserted{telle} VA est ``de type CM'' après extension finie du corps de base.

% lines 15–21
On est ramené, par ce qui précède, au cas d'un \deleted{VA} \inserted{schéma abélien} sur l'anneau des entiers \inserted{$R$} d'un corps de nbs $K$, avec action abélienne de $\operatorname{Gal}(\bar K/K)$ donc de $\pi_1(S,\unsure{\xi|\eta})$ ($S=\operatorname{Spec}R$). L'action de $\pi_1(S,\xi)$ \inserted{sur $E_\ell=T_\ell(A(\bar K))_{\mathbb Q_\ell}$} est loc algébrique \inserted{en tout $m\in\operatorname{Max}(R)$} (conditions locales sur l'action des $\pi_1$ des corps locaux, i.e. des groupes de décomposition) et est ``\inserted{loc} rationnelle'', \inserted{i.e. les pol. car. des \unsure{(v.b. de Frob)|subst. de Frob} à coeff $\in\mathbb Q$.} \deleted{\unsure{car}} Or on prouve que si on a deux telles représentations, \deleted{\unsure{alors}} un $E_\ell$ \uline{et}
